\subsection{The \acr{PHI} Instruction}

When control flow merges from multiple basic blocks, \acr{SSA} form requires a mechanism to select the appropriate value based on which predecessor block was executed. The \texttt{phi} instruction serves this purpose.

\begin{lstlisting}[style=ir]
define i32 @abs(i32 %x) {
entry:
  %is_neg = icmp slt i32 %x, 0
  br i1 %is_neg, label %negate, label %done

negate:
  %neg = sub i32 0, %x
  br label %done

done:
  %result = phi i32 [ %x, %entry ], [ %neg, %negate ]
  ret i32 %result
}
\end{lstlisting}

The \texttt{phi} instruction syntax is:

\begin{lstlisting}[style=ir]
%result = phi <type> [ <value1>, %block1 ], [ <value2>, %block2 ], ...
\end{lstlisting}

Semantics: "If control flow arrived from \texttt{\%block1}, use \texttt{<value1>}; if from \texttt{\%block2}, use \texttt{<value2>}."

\textbf{Why \texttt{phi} is necessary:}

Without \texttt{phi}, we cannot maintain \acr{SSA} form when values merge from different control paths. Consider this invalid attempt:

\begin{lstlisting}[style=ir]
; INVALID: Cannot assign to %result in multiple blocks
entry:
  %result = ... ; First assignment
  br ...

other_block:
  %result = ... ; ERROR: violates SSA (redefinition)
\end{lstlisting}

The \texttt{phi} instruction solves this by creating a single assignment point that selects from multiple incoming values.

\textbf{Multiple \texttt{phi} nodes:}

A block can have multiple \texttt{phi} instructions. All \texttt{phi} instructions in a block are conceptually executed simultaneously before any other instructions:

\begin{lstlisting}[style=ir]
loop:
  %i = phi i32 [ 0, %entry ], [ %i_next, %loop ]
  %sum = phi i32 [ 0, %entry ], [ %sum_next, %loop ]
  %i_next = add i32 %i, 1
  %sum_next = add i32 %sum, %i
  %cond = icmp slt i32 %i_next, 10
  br i1 %cond, label %loop, label %exit
\end{lstlisting}

This example implements a loop that sums integers from 0 to 9, demonstrating how \texttt{phi} enables iterative control flow in \acr{SSA} form.
